% Abkürzungsverzeichnis
\chapter*{Abkürzungsverzeichnis}
\vspace{-5em}

\addcontentsline{toc}{chapter}{Abkürzungsverzeichnis}
\begin{acronym}[LONGEST]
\acro{IPA}{Individuelle praktische Arbeit}
\acro{3-2-1-1-0}{Backup-Strategie (3 Kopien, 2 Medientypen, 1 externer Ort, 1 Offline-Kopie, 0 Fehler)}
\acro{API}{Application Programming Interface}
\acro{BPMN}{Business Process Model and Notation}
\acro{BV}{Berufliche Vorsorge}
\acro{CD}{Continuous Deployment}
\acro{CI/CD}{Continuous Integration / Continuous Deployment}
\acro{DN}{Distinguished Name}
\acro{E2E}{End-to-End}
\acro{EL}{Einzelleben}
\acro{Git}{Versionsverwaltungssystem (Distributed Version Control System)}
\acro{IDP}{Identity Provider}
\acro{IPA}{Individuelle Praktische Arbeit}
\acro{IST}{Aktueller Zustand (im Kontext: IST-Stand)}
\acro{KL}{Kollektivleben}
\acro{LDAP}{Lightweight Directory Access Protocol}
\acro{OpenLDAP}{Open Lightweight Directory Access Protocol}
\acro{OU}{Organizational Unit}
\acro{PoC}{Proof of Concept}
\acro{PV}{Private Vorsorge}
\acro{RxJS}{Reactive Extensions for JavaScript}
\acro{SOAP}{Simple Object Access Protocol}
\acro{SOLL}{Zielzustand (im Kontext: SOLL-Zeitplan)}
\acro{IST}{Aktueller Zustand (im Kontext: IST-Zeitplan)}
\acro{SSO}{Single Sign-On}
\acro{WAM}{Web Access Management}
\acro{XML}{Extensible Markup Language}
\acro{YAML}{Yet Another Markup Language}
\end{acronym} 